\section{Cadena de evidencia: ¿la Upstream Quality se traduce en Downstream Quality?}

Esta sección pasa del diseño del sistema a la evidencia. Si el sparse model solo obtiene una pretraining loss más baja gracias a un mayor número total de parámetros, la mejora en las tareas posteriores podría ser solo una consecuencia de la upstream quality y no de la forma de los parámetros en sí. Por ello, el expositor plantea una pregunta de control: con la misma upstream quality, ¿los sparse parameters adicionales siguen mejorando por sí mismos la downstream task?

\lecturefigure{slide-27-upstream-quality-research-question.jpg}{Pregunta de investigación: una vez fijada la upstream pretraining quality, ¿la forma del total de parámetros sigue influyendo en las tareas posteriores?}{Stanford Online, video oficial 00:43:58--00:44:17.}

\begin{knowledgebox}{Concepto previo: upstream y downstream}
La upstream task es el objetivo de preentrenamiento a gran escala, por ejemplo, el text-to-text denoising sobre C4; las downstream tasks son evaluaciones de transferencia como SuperGLUE o TriviaQA. Que ambas estén correlacionadas no significa que sean equivalentes: modelos con la misma pretraining loss pueden tener representaciones, distribuciones de capacidad y comportamientos de fine-tuning distintos.
\end{knowledgebox}

\subsection{SuperGLUE: la correlación es clara, pero no prueba la causalidad de los parámetros}

En la figura, la upstream quality de los modelos dense y Switch se correlaciona con la mejora en SuperGLUE. Si ambos tipos de puntos caen, en general, cerca de la misma tendencia, gran parte del beneficio downstream puede explicarse por una mejor pretraining quality; solo si los puntos sparse quedan sistemáticamente más arriba al mismo nivel de upstream se respalda que la forma de los parámetros aporta algo adicional.

\lecturefigure{slide-28-upstream-vs-superglue.jpg}{Relación entre la upstream quality y la SuperGLUE quality en Dense y Switch.}{Stanford Online, video oficial 00:44:18--00:45:19.}

\begin{knowledgebox}{Lectura de la figura: compara en horizontal, no solo a lo largo de la curva}
Primero compara los puntos dense y sparse en valores cercanos de upstream quality sobre el eje horizontal; después observa el error y la escala del modelo. Ver solamente que ambos suben conforme mejora el upstream demuestra una correlación, pero no prueba que «más parámetros dispersos sean eficaces por sí mismos a igual calidad».
\end{knowledgebox}

\begin{warningbox}{Correlation no es identificación del mecanismo}
Las familias de modelos cambian a la vez el número de parámetros, el tiempo de entrenamiento, la regularización, el enrutamiento y la estabilidad de la optimización. La tendencia de una gráfica de dispersión no basta para identificar cuál de estos factores provoca la diferencia downstream; se requieren experimentos con matched-quality, matched-compute y varias semillas aleatorias.
\end{warningbox}

Switch-C lleva el total de parámetros a cerca de 1.6T y aparece en la figura como un nuevo punto extremo. Su importancia está en mostrar que el sistema puede entrenarse y que el scaling continúa, no en anunciar que todas las tareas obtienen una mejora proporcional por la cifra de «un billón». Un punto extremo también puede ejercer un apalancamiento muy fuerte sobre la tendencia de la regresión, por lo que debe analizarse junto con los resultados de escala intermedia.

\lecturefigure{slide-29-switch-c-trillion-parameter.jpg}{Al añadir Switch-C, con cerca de 1.6T parámetros, se observa la tendencia de escalamiento de upstream y SuperGLUE.}{Stanford Online, video oficial 00:45:20--00:46:01.}

\teachervoice{El expositor no presenta Switch-C como el único atractivo, sino que lo usa para comprobar si la tendencia continúa. En términos de ingeniería, la pregunta más importante es si, con el mismo tiempo de dispositivo, la misma active compute y las mismas restricciones de servicio, vale más la pena que una configuración sparse o dense más pequeña.}

\subsection{TriviaQA: las tareas knowledge-heavy ofrecen otra perspectiva}

TriviaQA se orienta más a preguntas y respuestas sobre conocimiento, por lo que sirve para poner a prueba la intuición de que «más parámetros pueden almacenar conocimiento». En la figura, la upstream quality sigue correlacionada positivamente con la puntuación en TriviaQA, pero una sola gráfica de correlación no basta para atribuir la mejora a la capacidad de almacenamiento: la cobertura del corpus de entrenamiento, la tokenización, el fine-tuning y las reglas de coincidencia de respuestas influyen en el resultado.

\lecturefigure{slide-30-upstream-vs-triviaqa.jpg}{Relación entre la upstream quality y la calidad en TriviaQA.}{Stanford Online, video oficial 00:46:02--00:46:27.}

\begin{knowledgebox}{Lectura de la figura: qué respalda esta figura y qué no}
Respalda que «un mejor preentrenamiento suele ir acompañado de un mejor resultado en TriviaQA» y aporta evidencia de que vale la pena seguir investigando la sparse capacity; no prueba que el modelo almacene hechos internamente de forma localizable, ni que el número de parámetros importe más que la cobertura de datos, la recuperación o la active compute.
\end{knowledgebox}

\subsection{Fine-tuning de Base/Large y entrenamiento en 101 idiomas}

Esta sección traslada la evidencia del modelo más grande a escalas de uso más común y a distribuciones multilingües. Si el método sparse solo funcionara en el modelo más grande, su valor de ingeniería sería muy limitado. La clase muestra resultados de fine-tuning en los niveles Base y Large, lo que indica que también se obtienen beneficios a escalas más comunes; además, los experimentos multilingües usan mT5 como dense baseline y observan una ventaja general en 101 idiomas.

\lecturefigure{slide-31-fine-tuning-base-and-large.jpg}{Resultados de fine-tuning a escala Base y Large.}{Stanford Online, video oficial 00:46:28--00:47:07.}

\begin{knowledgebox}{Lectura de la tabla: además del promedio, busca las tareas que fallan}
Primero compara la calidad promedio y los FLOPs de dense y sparse a la misma escala; después revisa si alguna tarea se degrada. Una mejora promedio puede ocultar sensibilidad en tareas con pocos datos, secuencias largas o clases desbalanceadas; al desplegar, conviene conservar los resultados per-task en lugar de reportar solo el aggregate score.
\end{knowledgebox}

Un entorno multilingüe tiene de forma natural entradas heterogéneas, y el router puede formar una división de rutas por idioma, morfología o tema; pero eso no significa que cada idioma obtenga automáticamente un expert propio. Lo que realmente hay que revisar es la distribución entre idiomas de recursos altos y bajos, la cobertura del tokenizer, la expert utilization y si algunos idiomas se descartan sistemáticamente por overflow.

\lecturefigure{slide-32-multilingual-training.jpg}{Entrenamiento en 101 idiomas con mT5 como línea base y tendencia de la calidad.}{Stanford Online, video oficial 00:47:08--00:47:41.}

\teachervoice{El expositor presenta como un resultado importante que «los 101 idiomas mejoran», pero una reproducción en ingeniería debería reportar los resultados por idioma. Una mejora en el macropromedio no significa que ya estén resueltos los idiomas de pocos recursos, la equidad y la cobertura del enrutamiento.}

\begin{warningbox}{La división por idiomas puede generar una nueva cola larga}
Si el router concentra los idiomas de muchos recursos en experts estables, los idiomas de pocos recursos pueden recibir menos entrenamiento o un mayor dropping. Conviene medir por idioma la expert entropy, el overflow, el número de tokens y la downstream quality, para que el throughput global no oculte degradaciones locales.
\end{warningbox}

\subsection{Resumen de la sección}

La upstream quality es una variable explicativa importante de la downstream quality, pero una gráfica de correlación no prueba por sí sola el efecto causal independiente de los parámetros dispersos adicionales. Los resultados de Base/Large y multilingües amplían el ámbito de aplicación y, al mismo tiempo, exigen una auditoría más detallada por tarea y por idioma.
